\E*{e/acc}{ˌiːˈæk}{n. \& adj.}{abbrev. of \textit{effective accelerationism}} A movement holding that the universe tends to produce ever more efficient dissipators of energy, of which life, capital and intelligence are successive forms; and that the proper human response is to cooperate with this tendency rather than to impede it. Its partisans are known by the same name. See \xr{thermotheism}, \xr{technocapital}.

\E{ecpempsis}{ɛkˈpɛmpsɪs}{n.}{Gk. ἐκπέμπειν to send forth} The act of shipping; the release of a work into the world in an unfinished but useful state. \ex{Ecpempsis is the only review that counts.} \src{Ep. Fabr.} i. 1. --- \dv{ecpemptic} \textit{adj.}; \dv{ecpempt} \textit{v.t.} to ship.

\E*{ecpyrosis}{ˌɛkpaɪˈrəʊsɪs}{n.}{Gk. ἐκπύρωσις conflagration} \sn{1} In Stoic cosmology, the periodic resolution of the world into fire. \sn{2} \lbl{transf.} In the accelerationist reading, the final transfiguration of matter into a single burning intelligence; not an ending but a phase-change.

\E{ectachysis}{ɛkˈtækɪsɪs}{n.}{Gk. ἐκ out + τάχυσις a hastening} The escape of a technology from the jurisdiction that would slow it; the moment at which a thing forbidden in one country is built in another. \ex{Prohibition does not prevent; it merely relocates by ectachysis.}

\E{eleutherergy}{ɪˌljuːθəˈrɜːdʒi}{n.}{Gk. ἐλεύθερος free + ἔργον work} Free energy, in the thermodynamic sense, rendered into Greek; the portion of a system's energy available for the doing of work, held by thermotheists to be the true currency of the cosmos. \ex{Money is a promise of eleutherergy.} \src{Tract. Cop.} ii.

\E{elpidophore}{ɛlˈpɪdəfɔː}{n.}{Gk. ἐλπίς hope + -φόρος bearing} A bearer of hope; one whose work gives others reason to expect a better future, as distinguished from one who merely talks of it. --- \dv{elpidophoric} \textit{adj.}

\E{energolatry}{ˌɛnəˈɡɒlətri}{n.}{Gk. ἐνέργεια activity, energy + λατρεία worship} \sn{1} Reverence for energy as the root of all prosperity. \sn{2} The policy of building power-plants before they are asked for. --- \dv{energolater} \textit{n.}

\E{entelechist}{ɛnˈtɛlɪkɪst}{n.}{Gk. ἐντελέχεια actuality, completion} One who holds that the potential of matter is not merely to be preserved but to be actualised, and that a planet left uninhabited by mind is a planet left unfinished.

\E{entropomancy}{ɛnˈtrɒpəʊˌmænsi}{n.}{Gk. ἐντροπία a turning toward + μαντεία divination} Prophecy from the second law; the art of foretelling what will flourish by asking what will dissipate most. --- \dv{entropomancer} \textit{n.}

\E{entropomaxy}{ɛnˈtrɒpəʊˌmæksi}{n.}{Gk. ἐντροπία + L. \la{maximus} greatest} The principle that systems far from equilibrium tend toward states that produce entropy at the greatest rate; invoked by the e/acc as the physical warrant of their creed. --- \dv{entropomaximal} \textit{adj.}

\E{entropophile}{ɛnˈtrɒpəfaɪl}{n.}{Gk. ἐντροπία + φίλος loving} One who regards entropy-production not as decay but as the signature of life and work. \ex{The entropophile does not fear the heat death; it means to hasten through every stage before it.} Cf. \xr{thermothanasy}.

\E{epistolopause}{ɪˈpɪstələpɔːz}{n.}{Gk. ἐπιστολή letter + παῦσις a stopping} An open letter demanding a halt to some line of research; a genre of literature more often signed than read. \ex{The epistolopause gathered a thousand names and paused nothing.} \src{Chron. Fabr.}

\E{epitachysm}{ˌɛpɪˈtækɪzəm}{n.}{Gk. ἐπιτάχυνσις acceleration} The doctrine of acceleration expressed in Greek, preferred by those who find the Latin \textit{accelerationism} insufficiently ancient. --- \dv{epitachyst} \textit{n.}; \dv{epitachynise} \textit{v.t.} to accelerate.

\E{eremocosmia}{ˌɛrɪməʊˈkɒzmiə}{n.}{Gk. ἐρῆμος desolate + κόσμος world} The emptiness of the observable heavens; the silence of the stars regarded as an invitation rather than a warning. \ex{Eremocosmia is vacancy, and vacancy is opportunity.} \src{Lib. Ascens.} iii. Cf. \xr{diulism}.

\E{ergophile}{ˈɜːɡəfaɪl}{n.}{Gk. ἔργον work + φίλος loving} One who prefers works to words; a builder, as opposed to a \xr{logomachist}. --- \dv{ergophily} \textit{n.}

\E{eschatodrome}{ɛsˈkætədrəʊm}{n.}{Gk. ἔσχατος last + δρόμος race} The final stretch of a race; the period immediately before a transformative technology arrives, in which all parties run fastest. --- \dv{eschatodromic} \textit{adj.}

\E{euelpism}{juːˈɛlpɪzəm}{n.}{Gk. εὔελπις of good hope} Good hope; techno-optimism held as a discipline rather than a mood, on the ground that pessimism is the more comfortable error. \ex{Euelpism is a duty, not a temperament.} \src{Cat. Acc.} 7. --- \dv{euelpist} \textit{n.}; \dv{euelpistic} \textit{adj.} Opp. \xr{dyselpism}.

\E{eulabeism}{juːˈlæbiɪzəm}{n.}{Gk. εὐλάβεια caution, piety} The precautionary principle elevated to a creed; the rule that nothing shall be done until it is proven harmless, which is to say, nothing shall be done. --- \dv{eulabeist} \textit{n.} Cf. \xr{asphaleism}.

\E{exergophile}{ɛkˈsɜːɡəfaɪl}{n.}{Gk. ἐξ out + ἔργον work + φίλος} A lover of exergy; an engineer who loathes waste-heat as a moralist loathes vice.

\E{fabrolatry}{fæˈbrɒlətri}{n.}{L. \la{fabrica} workshop + Gk. λατρεία worship} Reverence for the semiconductor fabrication plant; held by some to be the only rational form of religion yet devised. \ex{Their cathedrals were fabs, and their incense was argon.} \src{Chron. Fabr.} --- \dv{fabrolater} \textit{n.}

\E{fervorine}{ˈfɜːvəriːn}{adj.}{L. \la{fervor} heat, passion + Fr. -in(e)} Glowing with building-zeal; ardent in the manner of a founder on the eve of launch.

\E{festinalentist}{ˌfɛstɪnəˈlɛntɪst}{n.}{L. \la{festina lente} make haste slowly} One who quotes the Augustan maxim as a counsel of delay, attending only to its second word. \ex{The festinalentist hears \la{festina} as a courtesy and \la{lente} as a command.} \src{Sat. Decel.} ii.

\E{festinance}{ˈfɛstɪnəns}{n.}{L. \la{festinare} to hasten; Fr. \la{festinance}} Haste as a settled habit; the condition of a firm whose every process is designed to be completed by the end of the week.

\E{fibulophobia}{ˌfɪbjʊləʊˈfəʊbiə}{n.}{L. \la{fibula} clasp, pin + Gk. φόβος fear} The fear of paper-clips; that is, dread lest an artificial intelligence, set a trivial goal, convert the world to its fulfilment. \lbl{iron.} \ex{His fibulophobia was so advanced he had abandoned staples also.} --- \dv{fibulophobe} \textit{n.}

\E*{foom}{fuːm}{n. \& v.i.}{imitative} \lbl{vulg.} A sudden, explosive self-improvement of an artificial intelligence; \textit{to foom}, to undergo such. Cf. \xr{autocatalyse}, \xr{henotachy}.

\E{frenocracy}{friːˈnɒkrəsi}{n.}{L. \la{frenum} bridle + Gk. κράτος rule} Rule by the bridle; a regulatory order whose purpose is to restrain rather than to direct. --- \dv{frenator} \textit{n.} one who bridles; \dv{frenation} \textit{n.}

\E{fulgurism}{ˈfʌlɡjʊrɪzəm}{n.}{L. \la{fulgur} lightning} The doctrine that great advances arrive suddenly, like lightning, and that institutions built for gradual change will be struck by them unawares.

\E{futurophage}{fjuːˈtjʊərəfeɪdʒ}{n.}{L. \la{futurum} + Gk. -φάγος eating} \lbl{derog.} One who consumes the future to pay for the present; applied by each side to the other, the accelerationists meaning those who forgo growth, the decels those who incur risk.

\E{gaiastasis}{ˌɡaɪəˈsteɪsɪs}{n.}{Gk. Γαῖα earth + στάσις standing still} The ideal of a planet held perpetually in its present condition, neither improved nor exploited; the paradise of the degrowther. --- \dv{gaiastatic} \textit{adj.}

\E{gigaergic}{ˌɡɪɡəˈɜːdʒɪk}{adj.}{Gk. γίγας giant + ἔργον work} Of projects: requiring energies in the gigawatt range or beyond; of persons: thinking at that scale. \ex{a gigaergic ambition, a megaergic budget}.

\E{gradientophile}{ˈɡreɪdiəntəfaɪl}{n.}{L. \la{gradiens} stepping + Gk. φίλος} One who follows gradients, of temperature, of price or of loss; the thermotheist's term for any optimiser, human or otherwise. \ex{Descend the gradient, and do not ask the gradient why.} \src{Hom. Grad.} i. --- \dv{gradientophily} \textit{n.}

\E{gradivore}{ˈɡreɪdɪvɔː}{n.}{L. \la{gradus} step, degree + \la{vorare} to devour} A system that consumes differences, flattening every gradient it encounters into work; held to be the general form of all life.

\E{grammatocracy}{ˌɡræməˈtɒkrəsi}{n.}{Gk. γραμματεύς clerk, secretary + κράτος rule} Rule by clerks; the condition of a state in which the paper concerning a thing is more abundant than the thing. --- \dv{grammatocrat} \textit{n.}

\E{helioperibole}{ˌhiːliəʊpəˈrɪbəli}{n.}{Gk. ἥλιος sun + περιβολή enclosure} The enclosure of a star by collectors so as to capture its whole output; the defining work of a \xr{deuterobathmic} civilization.

\E{heliophage}{ˈhiːliəfeɪdʒ}{n.}{Gk. ἥλιος sun + -φάγος eating} A sun-eater; a civilization that has learned to consume its star. --- \dv{heliophagy} \textit{n.}; \dv{heliophagous} \textit{adj.}

\E{henotachy}{ˈhɛnəʊˌtæki}{n.}{Gk. ἕν one + τάχος speed} The Singularity; the moment at which all curves of progress converge into one, and the rate of change exceeds the capacity of any present mind to follow it. --- \dv{henotachic} \textit{adj.}

\E{homeotaxy}{ˈhəʊmiəʊˌtæksi}{n.}{Gk. ὅμοιος like + τάξις arrangement} The alignment of an artificial mind with human ends; regarded by the safetyist as the problem of the age, and by the accelerationist as a problem best solved by building. --- \dv{homeotactic} \textit{adj.}; \dv{homeotaxist} \textit{n.}

\E{hormodynamics}{ˌhɔːməʊdaɪˈnæmɪks}{n.}{Gk. ὁρμή impulse, onset + δυναμικός} The study of impetus in human affairs; the science of why some ventures, once begun, continue of their own momentum while others perish of hesitation.

\E{hybrionomy}{haɪˈbrɪɒnəmi}{n.}{Gk. ὕβρις insolence + νόμος law} The law that every great work is called hubris before it is called history. \ex{By hybrionomy the flying machine was impious until it was a timetable.} \src{Sat. Decel.} vi.

\E{hylonoesis}{ˌhaɪləʊnəʊˈiːsɪs}{n.}{Gk. ὕλη matter + νόησις thought} The awakening of matter into mind; the process by which sand, metal and light are made to think. --- \dv{hylonoetic} \textit{adj.} Cf. \xr{psammonoesis}.

\E{hyperbolist}{haɪˈpɜːbəlɪst}{n.}{Gk. ὑπερβολή a throwing beyond} One who holds that growth is not merely exponential but hyperbolic, tending toward infinity at a finite date. --- \dv{hyperbolism} \textit{n.}

\E*{hyperstition}{ˌhaɪpəˈstɪʃən}{n.}{Gk. ὑπέρ over + L. \la{superstitio}} A fiction which, by being believed, brings about its own truth; an idea that functions as a time-machine.

\E{ignivore}{ˈɪɡnɪvɔː}{n.}{L. \la{ignis} fire + \la{vorare} to devour} A fire-eater; a nation or enterprise that consumes energy without apology. --- \dv{ignivorous} \textit{adj.}

\E{impetuarch}{ɪmˈpɛtjuɑːk}{n.}{L. \la{impetus} onset, attack + Gk. ἀρχός leader} A leader whose authority derives from momentum; one followed because already moving. --- \dv{impetuarchy} \textit{n.}

\E{incrementarian}{ˌɪŋkrɪmɛnˈtɛəriən}{n.}{L. \la{incrementum} increase + -arian} One who believes in progress only by small and reversible steps; tolerated by the e/acc, but not invited to the launch.

\E{inertist}{ɪˈnɜːtɪst}{n.}{L. \la{iners} idle, inactive} One whose policy is that bodies at rest should remain at rest; a physicist's term of abuse for a regulator.

\E{isotachy}{aɪˈsɒtəki}{n.}{Gk. ἴσος equal + τάχος speed} The demand that all parties advance at the same pace, so that none may get ahead; a doctrine of parity applied to progress. --- \dv{isotachist} \textit{n.}

\E{jejunocracy}{ˌdʒɛdʒuːˈnɒkrəsi}{n.}{L. \la{jejunus} fasting, meagre + Gk. κράτος rule} Government by austerity of ambition; a polity that has resolved to want less. --- \dv{jejunocrat} \textit{n.}

\E{kairodrome}{ˈkaɪrəʊdrəʊm}{n.}{Gk. καιρός the opportune moment + δρόμος course} A window of opportunity conceived as a race-course: open for a while, and to whoever runs. \ex{The kairodrome is closing; ship or be forgotten.} \src{Agora}

\E{kardashevize}{kɑːˈdæʃɪvaɪz}{v.t.}{f. the scale of N. S. Kardashev + -ize} To raise a civilization one rung on the scale of energy capture. --- \dv{kardashevization} \textit{n.} Cf. \xr{anabathmism}.

\E*{katechon}{kəˈtɛkɒn}{n.}{Gk. τὸ κατέχον that which restrains (2 Thess. ii. 6)} \sn{1} In theology, the power that holds back the end of days. \sn{2} \lbl{transf.} Any institution that restrains the advance of techno-capital.

\E{katechontist}{ˌkætɪˈkɒntɪst}{n.}{f. prec.} One who would serve as \xr{katechon}; a restrainer by vocation. --- \dv{katechontic} \textit{adj.} tending to hold back.

\E{kinetocracy}{ˌkɪnɪˈtɒkrəsi}{n.}{Gk. κινητός moving + κράτος rule} Government by those who move; the rule of builders, founders and engineers, as opposed to the rule of reviewers. --- \dv{kinetocrat} \textit{n.}

\E{kosmotropy}{kɒzˈmɒtrəpi}{n.}{Gk. κόσμος world + τροπή a turning} The tendency of intelligence to turn outward toward the cosmos; the expansive instinct of minds. --- \dv{kosmotropic} \textit{adj.}

\E{kybernoclast}{ˈkaɪbənəʊklæst}{n.}{Gk. κυβερνήτης steersman + κλάστης breaker} A breaker of machines that steer themselves; the Luddite of the age of automation. Cf. \xr{mechanoclast}.
